\subsubsection{模型建立}

针对问题三，从销量关系响应、相关情景生成和宏观结构判断三个方面建立模型。

\textbf{步骤1 关系响应}

设行归一化响应矩阵为$B^{(h)}$，弱、中、强档系数$h$分别取0.15、0.35和0.55。$B^{(h)}$ 的非对角元素刻画作物间销量响应的方向与强度：正值表示互补（一种作物销量增长带动另一种），负值表示替代（一种作物增长挤压另一种），行归一化保证响应幅度不随作物数放大。若问题二中心销量增长向量为$u_t$，则
\begin{equation}
r_t^{(h)}=u_t+B^{(h)}u_t,\qquad D_{tj}^{(h)}=D_{t-1,j}^{(h)}\left(1+r_{tj}^{(h)}\right).
\end{equation}
增长率超出题目区间时按边界裁剪。三档分别调用与问题二相同的MILP，中档方案作为正式代表。

\textbf{步骤2 相关情景生成}

销量增长、亩产量增长和随机食用菌价格增长的目标相关矩阵为
\begin{equation}
\Sigma^{(h)}=\begin{bmatrix}1&0.25h&0.55h\\0.25h&1&-0.20h\\0.55h&-0.20h&1\end{bmatrix}.
\end{equation}
对独立标准正态变量作Cholesky变换\upcite{golub2013}，再映射到问题二的有界边际区间。经验相关系数直接从最终导出的销量、亩产量和价格增长数组计算。

\textbf{步骤3 宏观结构判据}

两方案的41维累计面积结构相似度定义为
\begin{equation}
S_A(x,x')=\frac{\sum_{j=1}^{41}\min(A_j,A'_j)}{\sum_{j=1}^{41}\max(A_j,A'_j)}.
\end{equation}
当收益差小于1\%且$S_A\ge80\%$时，判定为宏观等价、微观平移。只有收益差大于1\%且$S_A<80\%$时才触发结构性备用方案。其中 $80\%$ 为本文用于辅助判断的操作性阈值，并非统计学或经济学意义上的客观临界值，其稳健性在模型检验中进一步讨论。

\subsubsection{模型求解}

中档关系方案在半价模式下平均收益为6887.03万元，比问题二基线高12.14万元；滞销压力下平均收益优势为62.41万元。正式200个情景中未出现模拟亏损。三档最大相关误差为0.0353，低于0.08阈值。需要说明的是，此处问题二基线在问题三的相关情景集下重新评价，其半价均值6874.89万元与问题二自身表格中的6891.51万元不同，差异仅来自情景生成方式，不影响配对比较。

半价模式下，弱、中、强三档平均收益分别为6909.83万元、6887.03万元和6865.49万元。随关系强度增强，收益标准差由约361万元增至534万元，5\%分位数由约6331万元降至5986万元，说明更强的模拟相关结构放大了收益波动（图\ref{fig:q3_strength}）。三档均使用24至25种作物，前三作物面积占比约44\%，未发生集中退化。
\begin{figure}[H]
  \centering
  \includegraphics[width=0.78\textwidth]{figures/q3_relationship_strength_sensitivity.png}
  \caption{问题三模拟关系强度的收益敏感性}
  \label{fig:q3_strength}
\end{figure}

三档方案两两收益差为0.061\%至0.186\%，面积结构相似度为94.73\%至96.63\%。图\ref{fig:q3_macro}中三组比较均落入微观等价平移区域。这一结果说明，在题面给定的土地适宜性、轮作和设施条件下，宏观种植结构由刚性生产约束主导，作物关系与相关结构的扰动主要体现为地块层面的微观调度，而非整体种植格局的改变。因此问题三的结论重点不在于「收益因关系结构而提高」，而在于「宏观结构稳定、微观调度灵活」——市场关系扰动在刚性生产约束面前只改变地块间的重新分配。
\begin{figure}[H]
  \centering
  \includegraphics[width=0.78\textwidth]{figures/q3_macro_structure_stability.png}
  \caption{问题三宏观面积结构稳定判据}
  \label{fig:q3_macro}
\end{figure}

\begin{longtable}{>{\centering\arraybackslash}p{0.13\textwidth}>{\centering\arraybackslash}p{0.10\textwidth}>{\centering\arraybackslash}p{0.15\textwidth}>{\centering\arraybackslash}p{0.15\textwidth}>{\centering\arraybackslash}p{0.15\textwidth}>{\centering\arraybackslash}p{0.13\textwidth}>{\centering\arraybackslash}p{0.08\textwidth}}
  \caption{问题三中档方案与问题二基线比较}\label{tab:q3_compare}\\
  \toprule
  方法 & 模式 & 均值/万元 & 5\%分位/万元 & 下尾均值/万元 & 最低值/万元 & 亏损比例 \\
  \midrule
  \endfirsthead
  \caption{问题三中档方案与问题二基线比较（续）}\\
  \toprule
  方法 & 模式 & 均值/万元 & 5\%分位/万元 & 下尾均值/万元 & 最低值/万元 & 亏损比例 \\
  \midrule
  \endhead
  \bottomrule
  \endfoot
  Q3中档 & 半价 & 6887.03 & 6125.76 & 6014.02 & 5732.33 & 0 \\
  Q2基线 & 半价 & 6874.89 & 6111.25 & 6009.93 & 5736.55 & 0 \\
  Q3中档 & 滞销 & 2625.40 & 2425.48 & 2371.49 & 2266.97 & 0 \\
  Q2基线 & 滞销 & 2562.99 & 2381.07 & 2328.86 & 2233.33 & 0 \\
\end{longtable}
